\documentclass[12pt,letterpaper]{article}
\renewcommand{\rmdefault}{phv}
\hbadness=99999

\usepackage[left=4cm,right=2cm,top=3cm,bottom=3cm]{geometry}
\usepackage[spanish,es-tabla]{babel}
\usepackage{fontenc}
\usepackage[utf8]{inputenc}
\usepackage{amsmath}
\usepackage{amssymb,amsfonts,cancel}
\usepackage{array}
\usepackage{enumerate}
\usepackage{etoolbox}
\usepackage{fancyhdr}
\usepackage{float}
\usepackage{graphicx}
\usepackage[pdfstartpage=1,backref=false,pagebackref=false,pdfborder={0 0 0},colorlinks=true,linkcolor=blue,citecolor=blue,urlcolor=blue]{hyperref}
\usepackage{longtable}
\usepackage{multicol}
\usepackage{multirow}
\usepackage{ragged2e}
\usepackage{caption}
\usepackage{subcaption}
\usepackage{titlesec}
\usepackage{titletoc}
\usepackage{titling}
\usepackage{colortbl}
\usepackage{xcolor}
\usepackage{tabularx}
\usepackage{listings}
\usepackage{booktabs}
\usepackage[round]{natbib}

% Configuración de encabezados y pies de página
\fancyhead{}
\fancyfoot[c]{\thepage}
\pagestyle{fancy}

% Configuración de tablas de contenido
\usepackage{tocloft}
\addto\captionsspanish{\renewcommand{\contentsname}{Contenido}}
\renewcommand{\cfttoctitlefont}{\hfill\normalsize}
\renewcommand{\cftaftertoctitle}{\hfill}
\renewcommand{\cftloftitlefont}{\hfill\normalsize}
\renewcommand{\cftafterloftitle}{\hfill\mbox{}}
\renewcommand{\cftlottitlefont}{\hfill\normalsize}
\renewcommand{\cftafterlottitle}{\hfill\mbox{}}

\renewcommand{\cftfigfont}{Fig. }
\renewcommand{\cfttabfont}{TABLA }
\renewcommand{\tablename}{TABLA}

\usepackage{layout}

% Formato de títulos
\titleclass{\subsubsubsection}{straight}[\subsection]
\newcounter{subsubsubsection}[subsubsection]
\renewcommand\thesubsubsubsection{\thesubsubsection.\arabic{subsubsubsection}}
\renewcommand\theparagraph{\thesubsubsubsection.\arabic{paragraph}}

\titleformat{\section}[block]{\bfseries\centering}{\thesection.}{1em}{}
\titleformat{\subsection}
{\normalfont\normalsize\bfseries\MakeUppercase}{\thesubsection}{1em}{}
\titlespacing*{\subsection}{0pt}{0em}{0em}
\titleformat{\subsubsection}
{\normalfont\normalsize\bfseries}{\thesubsubsection}{1em}{}
\titlespacing*{\subsubsection}{0pt}{0em}{0em}
\titleformat{\subsubsubsection}
{\normalfont\normalsize\bfseries}{\thesubsubsubsection}{1em}{}
\titlespacing*{\subsubsubsection}{0pt}{0em}{0em}

\makeatletter
\renewcommand\paragraph{\@startsection{paragraph}{5}{\z@}%
{0em}%
{-1em}%
{\normalfont\normalsize\bfseries}}
\renewcommand\subparagraph{\@startsection{subparagraph}{6}{\parindent}%
{0em}%
{-1em}%
{\normalfont\normalsize\bfseries}}
\def\toclevel@subsection{2}
\def\toclevel@subsubsection{3}
\def\toclevel@subsubsubsection{4}
\def\toclevel@paragraph{5}
\def\toclevel@subparagraph{6}
\setlength{\cftsecindent}{0pt}
\setlength{\cftsubsecindent}{0pt}
\setlength{\cftsubsubsecindent}{0pt}
\setlength{\cftfigindent}{0pt}
\setlength{\cfttabindent}{0pt}

\g@addto@macro\cftsecpagefont{\bfseries}
\g@addto@macro\cftsubsecpagefont{\bfseries\MakeUppercase}
\g@addto@macro\cftparapagefont{\bfseries}
\g@addto@macro\cftsubparapagefont{\bfseries}
\g@addto@macro\cftfigpagefont{\bfseries}
\g@addto@macro\cfttabpagefont{\bfseries}
\makeatother

\setcounter{secnumdepth}{4}
\setcounter{tocdepth}{4}

\makeatletter
\newcolumntype{U}{>{$}c<{$}}
\renewcommand{\fnum@table}{\textbf{\tablename~\thetable}.
\sffamily}
\renewcommand\thetable{\Roman{table}}
\makeatother

\renewcommand{\headrulewidth}{0pt}
\raggedbottom

\setlength{\arrayrulewidth}{1pt}
\setlength{\doublerulesep}{0mm}
\setlength{\parindent}{0em}
\setlength{\parskip}{2em}

% Colores y Listings
\definecolor{codegreen}{rgb}{0,0.6,0}
\definecolor{codegray}{rgb}{0.5,0.5,0.5}
\definecolor{codepurple}{rgb}{0.58,0,0.82}
\definecolor{backcolour}{rgb}{0.95,0.95,0.92}

\lstdefinestyle{pythonstyle}{
    backgroundcolor=\color{backcolour},
    commentstyle=\color{codegreen},
    keywordstyle=\color{magenta},
    numberstyle=\tiny\color{codegray},
    stringstyle=\color{codepurple},
    basicstyle=\ttfamily\footnotesize,
    breakatwhitespace=false,
    breaklines=true,
    captionpos=b,
    keepspaces=true,
    numbers=left,
    numbersep=5pt,
    showspaces=false,
    showstringspaces=false,
    showtabs=false,
    tabsize=2,
    language=Python,
    literate={á}{{\'a}}1 {é}{{\'e}}1 {í}{{\'i}}1 {ó}{{\'o}}1 {ú}{{\'u}}1 {Á}{{\'A}}1 {É}{{\'E}}1 {Í}{{\'I}}1 {Ó}{{\'O}}1 {Ú}{{\'U}}1 {ñ}{{\~n}}1 {Ñ}{{\~N}}1
}
\lstset{style=pythonstyle}

\usepackage{tikz}
\usetikzlibrary{arrows.meta, positioning}

\title{Bitácora Metodológica de Hallazgos}
\author{Frank Daza}

\begin{document}
\maketitle
\tableofcontents

\section{Propósito y convenciones}

Este documento es el \textbf{registro empírico} de la ejecución del proyecto de grado.
A diferencia del anteproyecto en \texttt{docs/proyecto\_de\_grado/}, que es prospectivo,
la bitácora acumula evidencia conforme se completan las tareas del backlog.
Cada hallazgo queda ligado a artefactos en \texttt{results/} y \texttt{models/}.

\subsection{Flujo documental}

El anteproyecto define el problema y el diseño metodológico.
Esta bitácora registra lo ejecutado (tareas 3--16).
El documento final (\texttt{Trabajo de Grado - Frank Daza.tex}, tarea 17) sintetiza
los hallazgos con referencias cruzadas \texttt{\textbackslash ref\{hallazgo:task-N\}}.
Derivados posteriores: artículo científico (tarea 18) y ponencia (tarea 19).

\subsection{Trazabilidad obligatoria}

Toda tarea del backlog que produzca evidencia experimental debe registrar un hallazgo
con la etiqueta \texttt{\textbackslash label\{hallazgo:task-N\}}, donde \texttt{N}
es el identificador numérico de la tarea.
Una tarea sin esta etiqueta se considera \textbf{no documentada}, aunque haya generado
CSV en \texttt{results/}.

\subsection{Política de bibliografía}

La bibliografía es \textbf{compartida} con el anteproyecto:
\texttt{docs/proyecto\_de\_grado/Referencias.bib}.
Está prohibido duplicar entradas en un \texttt{.bib} local.
Las citas nuevas se agregan con el skill \texttt{agregar-cita} sobre el archivo compartido.
Motor: \texttt{natbib} + estilo \texttt{apalike}.

\subsection{Plantilla estándar de hallazgo}

Toda tarea posterior debe replicar la siguiente estructura en el fragmento
\texttt{hallazgos/} que corresponda a su fase:

\begin{verbatim}
\subsubsection{Task-N: Título de la tarea}
\label{hallazgo:task-N}
\begin{tabularx}{\textwidth}{|l|X|}\hline
\textbf{Fecha} & YYYY-MM-DD \\\hline
\textbf{Actividad} & A# del anteproyecto \\\hline
\textbf{Artefactos} & \texttt{results/...}, \texttt{models/...} \\\hline
\end{tabularx}

\paragraph{Objetivo.} Qué se buscaba verificar.
\paragraph{Procedimiento.} Resumen reproducible (semilla, hiperparámetros,
comando \texttt{uv run}).
\paragraph{Resultados.} Tablas y figuras con \texttt{\textbackslash input}
desde \texttt{results/figures/}.
\paragraph{Interpretación.} Qué implica para la hipótesis de eficiencia de datos.
\paragraph{Decisiones.} Cambios congelados (p.\,ej.\ $L=4$) o riesgos abiertos.
\paragraph{Limitaciones.} Sesgos, supuestos, deuda técnica.
\end{verbatim}

\subsection{Mapeo de fragmentos por fase}

\begin{tabularx}{\textwidth}{|l|X|}\hline
\textbf{Fragmento} & \textbf{Tareas} \\\hline
\texttt{h0\_fundamentos.tex} & 1, 2, 3, 4, 5, 6 \\\hline
\texttt{h1\_arquitectura.tex} & 7, 8, 9, 10, 11 \\\hline
\texttt{h2\_experimentacion.tex} & 12, 13 \\\hline
\texttt{h3\_analisis.tex} & 14, 15, 16 \\\hline
\end{tabularx}

\section{Fase 0 — Fundamentos reproducibles y datos}
% Fragmento: Fase 0 — Fundamentos reproducibles y datos
% Tareas asignadas: TASK-1, TASK-2, TASK-3, TASK-4, TASK-5, TASK-6
% Plantilla: ver sección «Propósito y convenciones» de la bitácora.
% Obligatorio: \label{hallazgo:task-N} por cada tarea documentada.

\subsubsection{Task-1: Infraestructura y entorno reproducible}
\label{hallazgo:task-1}
\begin{tabularx}{\textwidth}{|l|X|}\hline
\textbf{Fecha} & 2026-08-16 \\\hline
\textbf{Actividad} & Infraestructura (m-0) \\\hline
\textbf{Artefactos} & \texttt{src/config.py}, \texttt{src/utils/seed.py}, \texttt{src/utils/device.py}, \texttt{src/smoke.py}, \texttt{uv.lock}, \texttt{.python-version} \\\hline
\end{tabularx}

\paragraph{Objetivo.}
Establecer un entorno reproducible con una única fuente de hiperparámetros,
semilla global fijada y selección de dispositivo \texttt{cuda > mps > cpu}.

\paragraph{Procedimiento.}
Se implementó \texttt{ExperimentConfig} como dataclass congelada, \texttt{set\_seed(42)}
con \texttt{torch.use\_deterministic\_algorithms(True, warn\_only=True)} y
\texttt{get\_device()} con la prelación definida en \texttt{AGENTS.md}.
Verificación con:
\begin{verbatim}
uv run python -m src.smoke
\end{verbatim}

\paragraph{Resultados.}
Stack congelado: Python 3.12, PyTorch 2.9.1, torchvision 0.24.1,
PennyLane 0.45.1, NumPy $\geq 2.0,< 2.3$.
Dispositivo detectado: \textbf{mps}.
Pérdida de referencia (dos corridas idénticas): \textbf{1.54205954}.
Resultado reproducibilidad: \textbf{PASS}.

\paragraph{Interpretación.}
El entorno local cumple el criterio de reproducibilidad bit a bit dentro del mismo
dispositivo. La semilla 42 queda fijada por convención del anteproyecto.

\paragraph{Decisiones.}
Semilla global 42 en \texttt{ExperimentConfig.semilla}.
Gestor de paquetes exclusivo: UV (\texttt{uv sync --frozen}).
Dispositivo registrado en cada corrida (contrato de tarea 4).

\paragraph{Limitaciones.}
La reproducibilidad bit a bit solo se garantiza dentro del mismo dispositivo;
en \texttt{mps} algunas operaciones caen a CPU sin equivalente completo de
\texttt{torch.cuda.manual\_seed\_all}.

\subsubsection{Task-2: Estructura LaTeX del trabajo de grado y bitácora de hallazgos}
\label{hallazgo:task-2}
\begin{tabularx}{\textwidth}{|l|X|}\hline
\textbf{Fecha} & 2026-08-16 \\\hline
\textbf{Actividad} & Documentación (m-0) \\\hline
\textbf{Artefactos} & \texttt{docs/trabajo\_de\_grado/} (bitácora, preámbulo, esqueleto de tesis, fragmentos \texttt{hallazgos/}) \\\hline
\end{tabularx}

\paragraph{Objetivo.}
Crear la estructura LaTeX del documento final y la bitácora metodológica de hallazgos,
con trazabilidad \texttt{task-N} $\leftrightarrow$ artefactos en \texttt{results/}.

\paragraph{Procedimiento.}
Se copió el preámbulo UAO del anteproyecto a \texttt{preambulo\_uao.tex} sin modificar
el original. Se crearon los fragmentos \texttt{h0\_fundamentos.tex} a \texttt{h3\_analisis.tex},
el esqueleto \texttt{Trabajo de Grado - Frank Daza.tex} con capítulos vacíos y la cadena
de compilación \texttt{pdflatex} $\rightarrow$ \texttt{bibtex} $\rightarrow$ \texttt{pdflatex}
contra \texttt{../proyecto\_de\_grado/Referencias.bib}.

\paragraph{Resultados.}
Estructura creada en \texttt{docs/trabajo\_de\_grado/}:
\texttt{preambulo\_uao.tex}, bitácora, esqueleto de tesis, carpetas \texttt{hallazgos/},
\texttt{capitulos/} y \texttt{Figuras/} (logo UAO).
Bibliografía compartida sin entradas duplicadas.
Cita de verificación: PennyLane~\cite{bergholm2018pennylane}.

\paragraph{Interpretación.}
A partir de esta tarea, cada hallazgo experimental se registra en el fragmento de fase
correspondiente. El documento final (tarea 17) sintetizará esta evidencia con
\texttt{\textbackslash ref\{hallazgo:task-N\}} en lugar de reconstruir resultados desde CSV.

\paragraph{Decisiones.}
Preámbulo copiado (no refactorizado) del anteproyecto entregado.
Fragmentos sin espacios en el nombre de archivo.
Logo UAO copiado (no symlink) para portabilidad entre máquinas.

\paragraph{Limitaciones.}
Los documentos principales (\texttt{Bitacora Metodologica de Hallazgos.tex},
\texttt{Trabajo de Grado - Frank Daza.tex}) llevan espacios en el nombre;
la compilación requiere comillas en la línea de comandos.

\subsubsection{Task-3: Auditoría del Brain Tumor MRI Dataset (A4)}
\label{hallazgo:task-3}
\begin{tabularx}{\textwidth}{|l|X|}\hline
\textbf{Fecha} & 2026-08-17 \\\hline
\textbf{Actividad} & A4 --- Curaduría y auditoría del dataset (m-0) \\\hline
\textbf{Artefactos} & \texttt{src/data/audit.py}, \texttt{results/dataset\_manifest.csv}, \texttt{results/dataset\_audit.md}, \texttt{results/figures/distribucion\_clases.png} \\\hline
\end{tabularx}

\paragraph{Objetivo.}
Verificar el conteo real, el balance por clase y partición, la integridad de las
imágenes y la presencia de duplicados exactos antes de diseñar escenarios de
escasez, siguiendo buenas prácticas de preparación de datos en imagen médica~\cite{willemink2020preparing}.

\paragraph{Procedimiento.}
Dataset local en \texttt{notebooks/data/}, enlazado por symlink a
\texttt{data/brain\_tumor\_mri/} (ruta de \texttt{ExperimentConfig}).
Recorrido estable con \texttt{pathlib}, verificación PIL (\texttt{verify} +
\texttt{load}), hash SHA-256 de contenido y clasificación de duplicados en tres
categorías. Ejecución:
\begin{verbatim}
uv run python -m src.data.audit
\end{verbatim}

\paragraph{Resultados.}
Total inventariado: \textbf{7023} imágenes (coincide con el anteproyecto).
Extensiones: solo \texttt{.jpg}. Imágenes corruptas: \textbf{0}.
Duplicados exactos (grupos con $n>1$): intra clase \textbf{194}, entre clases
\textbf{0}, Training$\leftrightarrow$Testing \textbf{79}.
Exclusiones: \textbf{297} filas (\texttt{duplicado\_exacto}: 194;
\texttt{fuga\_train\_test}: 103). Utilizables: \textbf{6726}.

\begin{center}
\begin{tabular}{lrrrr}\hline
\textbf{Clase} & \textbf{Testing} & \textbf{Training} & \textbf{Total} \\\hline
glioma       & 300 & 1321 & 1621 \\
meningioma   & 306 & 1339 & 1645 \\
notumor      & 405 & 1595 & 2000 \\
pituitary    & 300 & 1457 & 1757 \\\hline
\textbf{Total} & 1311 & 5712 & 7023 \\\hline
\end{tabular}
\end{center}

Heterogeneidad: modos \texttt{L} (3093), \texttt{RGB} (3926), \texttt{P}/\texttt{RGBA} (4);
resolución dominante $512\times512$ (4742 imágenes). Las cuatro clases corresponden
a categorías del sistema nervioso central en la clasificación WHO 2021~\cite{louis20212021}.

\begin{figure}[ht]
\centering
\includegraphics[width=0.85\textwidth]{../../../results/figures/distribucion_clases}
\caption{Distribución de clases por partición (imágenes no excluidas).}
\label{fig:distribucion-clases-task3}
\end{figure}

\paragraph{Interpretación.}
El dataset agregado de Nickparvar~\cite{nickparvar2023dataset} es heterogéneo en
modo de color y resolución; la auditoría confirma que el tamaño declarado es correcto
pero existen duplicados exactos y fuga potencial entre \texttt{Training/} y
\texttt{Testing/}. Sin esta verificación, la exactitud de modelos posteriores
podría inflarse por memorización de copias idénticas~\cite{haddou2025hqcm}.

\paragraph{Decisiones.}
Unir \texttt{Training/} y \texttt{Testing/} al conjunto completo para validación
cruzada estratificada $k=5$ (decisión D1, TASK-6); no se usa \texttt{Testing/} como
holdout externo. Las exclusiones quedan en \texttt{dataset\_manifest.csv} mediante
columnas \texttt{excluida} y \texttt{motivo\_exclusion}.

\paragraph{Limitaciones.}
El hash SHA-256 detecta duplicados byte a byte, no re-codificaciones JPEG con distinta
compresión. La conversión a 3 canales y el redimensionado a $224\times224$ quedan
para TASK-5.

% --- Tareas pendientes en este fragmento ---
% TASK-4: Contrato de métricas, logging y artefactos
% TASK-5: Pipeline de preprocesamiento y aumento de datos (A5)
% TASK-6: Protocolo de particiones: holdout, k-fold y fracciones (A7)


\section{Fase 1 — Arquitectura híbrida CNN-VQC}
% Fragmento: Fase 1 — Arquitectura híbrida CNN-VQC
% Tareas asignadas: TASK-7, TASK-8, TASK-9, TASK-10, TASK-11
% Plantilla: ver sección «Propósito y convenciones» de la bitácora.
% Obligatorio: \label{hallazgo:task-N} por cada tarea documentada.

% --- Tareas pendientes en este fragmento ---
% TASK-7: Extractor congelado y fábrica de backbones (A1)
% TASK-8: Bucle unificado de entrenamiento y evaluación (Trainer)
% TASK-9: Clasificador cuántico variacional de 4 qubits (A2)
% TASK-10: Modelo híbrido HQCNN con TorchLayer (A3)
% TASK-11: Ablación de profundidad L y presupuesto de cómputo


\section{Fases 2 y 3 — Experimentación y análisis}
% Fragmento: Fases 2 — Experimentación
% Tareas asignadas: TASK-12, TASK-13
% Plantilla: ver sección «Propósito y convenciones» de la bitácora.
% Obligatorio: \label{hallazgo:task-N} por cada tarea documentada.

% --- Tareas pendientes en este fragmento ---
% TASK-13: Campaña experimental k-fold en escenarios de escasez (A8)

\subsubsection{Task-12: Líneas base clásicas con el 100\% del conjunto (A6)}
\label{hallazgo:task-12}
\begin{tabularx}{\textwidth}{|l|X|}\hline
\textbf{Fecha} & 2026-08-17 \\\hline
\textbf{Actividad} & A6 --- Entrenamiento de las líneas base clásicas (m-2) \\\hline
\textbf{Artefactos} & \texttt{src/models/baseline.py}, \texttt{src/models/factory.py}, \texttt{src/experiments/baselines.py}, \texttt{tests/test\_baseline.py}, \texttt{results/experiments.csv}, \texttt{results/baselines\_summary.csv}, \texttt{models/*.pt}, \texttt{results/history/*.json} \\\hline
\end{tabularx}

\paragraph{Objetivo.}
Establecer el límite superior de exactitud contra el cual se interpretará el HQCNN y el resto
de la campaña factorial (TASK-13). Las líneas base EfficientNet-B0 y ResNet-50 comparten el
mismo \texttt{Trainer} (TASK-8), las mismas particiones de \texttt{results/splits.json}
(TASK-6) y el mismo presupuesto de 15 épocas, con extractor ImageNet congelado y cabeza
\texttt{Linear(dim $\rightarrow$ 4)} (TASK-7).

\paragraph{Procedimiento.}
\texttt{ClassicalBaseline} envuelve \texttt{build\_backbone} + \texttt{CabeceraReduccion} y
mantiene el backbone en \texttt{eval()} durante el entrenamiento. La fábrica
\texttt{build\_model(nombre, cfg)} unifica la construcción para baselines y HQCNN (sin ramas
\texttt{if} en el orquestador). Versiones de pesos: EfficientNet-B0
\texttt{IMAGENET1K\_V1}~\cite{tan2019efficientnet}; ResNet-50 \texttt{IMAGENET1K\_V2}~\cite{he2016deep}
(registradas en \texttt{VERSIONES\_PESOS}). Dispositivo de ejecución: \texttt{mps} (macOS).
Antes de cada corrida, \texttt{verificar\_indices\_fold} comprueba que los índices cargados
coinciden con \texttt{splits.json}. Ejecución:
\begin{verbatim}
uv run python -m src.experiments.baselines
uv run pytest tests/test_baseline.py -q
\end{verbatim}

\paragraph{Resultados por fold (validación, fracción 1.00).}
\begin{center}
\small
\begin{tabular}{llrrrr}\hline
\textbf{Modelo} & \textbf{Fold} & \textbf{Acc. val} & \textbf{F1 macro} & \textbf{F1 pond.} & \textbf{Tiempo (s)} \\\hline
EfficientNet-B0 & 0 & 0.9376 & 0.9368 & 0.9378 & 640 \\
EfficientNet-B0 & 1 & 0.9390 & 0.9377 & 0.9387 & 643 \\
EfficientNet-B0 & 2 & 0.9346 & 0.9335 & 0.9347 & 645 \\
EfficientNet-B0 & 3 & 0.9368 & 0.9360 & 0.9371 & 663 \\
EfficientNet-B0 & 4 & 0.9368 & 0.9356 & 0.9368 & 782 \\
ResNet-50 & 0 & 0.9190 & 0.9168 & 0.9181 & 1028 \\
ResNet-50 & 1 & 0.9152 & 0.9133 & 0.9148 & 982 \\
ResNet-50 & 2 & 0.9257 & 0.9240 & 0.9252 & 987 \\
ResNet-50 & 3 & 0.9182 & 0.9164 & 0.9177 & 1012 \\
ResNet-50 & 4 & 0.9234 & 0.9218 & 0.9231 & 942 \\\hline
\end{tabular}
\end{center}

\paragraph{Agregación sobre 5 folds (media $\pm$ desv. estándar).}
\begin{center}
\small
\begin{tabular}{lrrr}\hline
\textbf{Métrica} & \textbf{EfficientNet-B0} & \textbf{ResNet-50} \\\hline
Exactitud val. & $0.9369 \pm 0.0012$ & $0.9203 \pm 0.0039$ \\
F1 macro & $0.9358 \pm 0.0012$ & $0.9185 \pm 0.0040$ \\
F1 ponderado & $0.9370 \pm 0.0011$ & $0.9198 \pm 0.0039$ \\
Brecha $G$ & $0.0273 \pm 0.0021$ & $0.0229 \pm 0.0037$ \\
Inferencia (ms/lote) & $52.7 \pm 14.8$ & $140.9 \pm 73.0$ \\\hline
\end{tabular}
\end{center}

Pruebas automatizadas: \textbf{19 passed} (\texttt{test\_baseline.py} + \texttt{test\_backbones.py}).
Diez celdas en \texttt{experiments.csv} con clave \texttt{(modelo, 1.00, fold, semilla=42)}.

\paragraph{Contraste con la literatura.}
\citet{haddou2025hqcm} reportan 96.48\,\% de precisión con HQCM-EBTC sobre el mismo
\textit{dataset}, entrenando una CNN personalizada desde cero con el conjunto completo y el
\textit{split} original \texttt{Training/Testing}. \citet{khan2024brain} evalúan un HQCNN
híbrido en condiciones distintas (arquitectura, preprocesamiento y partición).
\textbf{Los números no son directamente comparables} con los de esta tarea porque:
(i)~validación cruzada estratificada de 5 folds frente a un único \textit{holdout};
(ii)~extractor preentrenado congelado frente a entrenamiento completo del backbone;
(iii)~cuello de botella a 4 dimensiones (restricción del diseño híbrido, ver
Hallazgo~\ref{hallazgo:task-7}). La exactitud observada ($\sim$93.7\,\% EfficientNet-B0,
$\sim$92.0\,\% ResNet-50) no indica fallo de implementación: es coherente con la restricción
de entrenar solo la cabeza lineal sobre características fijas.

\paragraph{Interpretación.}
EfficientNet-B0 supera a ResNet-50 bajo la misma restricción de extractor congelado y la
misma cabeza, con menor varianza entre folds (desv.\ std.\ de exactitud 0.12 vs.\ 0.39 puntos
porcentuales). La brecha de generalización $G$ se mantiene baja ($\sim$2--3 puntos),
coherente con pocos parámetros entrenables. Estas 10 celdas quedan disponibles para
reutilización por TASK-13 sin reentrenamiento (\texttt{corrida\_completada}).

\paragraph{Decisiones.}
ResNet-50 con \texttt{IMAGENET1K\_V2} (no intercambiable con V1 en comparaciones
bibliográficas). Dispositivo \texttt{mps} para baselines clásicas (a diferencia del HQCNN,
que usa CPU por inestabilidad de \texttt{TorchLayer} en MPS, Hallazgo~\ref{hallazgo:task-11}).

\paragraph{Limitaciones.}
El extractor congelado limita el techo alcanzable: la tabla \textbf{no} representa el máximo
de EfficientNet-B0 o ResNet-50 con \textit{fine-tuning} completo. Los 5 folds comparten
observaciones (TASK-6); la desviación estándar entre folds mide estabilidad de la partición,
no error de muestreo independiente. El tiempo de inferencia de ResNet-50 es mayor por la
mayor dimensión latente (2048 vs.\ 1280) y el costo del \textit{forward} congelado.

% Fragmento: Fase 3 — Análisis
% Tareas asignadas: TASK-14, TASK-15, TASK-16
% Plantilla: ver sección «Propósito y convenciones» de la bitácora.
% Obligatorio: \label{hallazgo:task-N} por cada tarea documentada.

% --- Tareas pendientes en este fragmento ---
% TASK-14: Consolidación multimétrica y de tiempos (A9)
% TASK-15: Análisis estadístico ANOVA y comparaciones post-hoc (A10)
% TASK-16: Curvas de aprendizaje y brecha de generalización (A11)


\section{Síntesis de decisiones metodológicas}

Esta sección consolida las decisiones congeladas a medida que avanza el proyecto:

\begin{itemize}
    \item \textbf{D1} --- Orden de particiones y fracciones de escasez (tarea 6).
    Ver~\ref{hallazgo:task-6} y Figura~\ref{fig:diagrama-d1-task6}.
    \begin{itemize}
        \item Unir \texttt{Training/} y \texttt{Testing/} tras auditoría (6726 imágenes utilizables).
        \item \texttt{StratifiedKFold} con $k=5$, \texttt{random\_state=42}, sobre el conjunto completo.
        \item Validación de tamaño fijo por fold; idéntica en las cuatro fracciones de escasez.
        \item Submuestreo estratificado \textbf{anidado} solo del entrenamiento: 10\,\% $\subset$ 25\,\% $\subset$ 50\,\% $\subset$ 100\,\%.
        \item Índices persistidos en \texttt{results/splits.json} con validación de hash del manifiesto.
        \item Holdout externo sobre \texttt{Testing/} descartado (duplicados cruzados y diseño del anteproyecto).
        \item Lectura literal de A8 (submuestrear antes de partir) rechazada: varianza de validación dominaría al 10\,\%.
    \end{itemize}
    \item \textbf{D2} --- Profundidad $L$ del VQC y presupuesto de cómputo (tarea 11).
    Ver~\ref{hallazgo:task-11} y Figura~\ref{fig:diagrama-d2-task11}.
    \begin{itemize}
        \item Ablación económica: $L \in \{2,4,6\}$, fold 0, fracción 25\,\%, 5 épocas (no A8).
        \item Criterio pre-declarado: mayor F1 macro; empate $\leq 0{,}02$ favorece $L$ menor.
        \item $L$ congelada en \texttt{results/selected\_hparams.json} (valor y artefactos en TASK-11).
        \item Extrapolación: horas HQCNN solo vs.\ cota conservadora ($3$ modelos); umbral 72 h.
        \item \textit{No-go} en CPU local si supera umbral; \textit{go} condicionado a sonda de 1 época
              en Colab/CUDA (TASK-13) antes de las 60 celdas.
        \item Sin recorte de épocas sin evidencia de convergencia; D2 (HQCNN al 100\,\%) se confirma tras re-extrapolar.
        \item Historiales por $L$ en \texttt{results/history/} con sufijo \texttt{\_L\{n\}}.
    \end{itemize}
\end{itemize}

\bibliographystyle{apalike}
\bibliography{../proyecto_de_grado/Referencias}
\end{document}
